\section{Постановка модели}

Полное состояние HANK-экономики обозначается через \(X_t^{HANK}\). Оно включает агрегированные переменные, распределение домохозяйств и экзогенные шоки. Это состояние является экономической мотивацией, но не используется напрямую как вход правила политики. В вычислительной части работы рассматривается редуцированное состояние
\[
x_t = R(X_t^{HANK}),
\]
которое состоит из агрегатного блока и распределительного блока:
\[
x_t =
\left(
\pi_t^{gap},
y_t^{gap},
r_t^*,
\overline{\kappa}_t,
\ell_t
\right).
\]

Здесь \(\overline{\kappa}_t\) обозначает среднюю предельную склонность к потреблению, а \(\ell_t\) -- долю низколиквидных домохозяйств.

Динамика состояния задаётся локальной системой
\[
x_{t+1}
=
A x_t
+
g(x_t,i_t)
+
\varepsilon_{t+1}.
\]
Матрица \(A\) описывает собственную динамику агрегатных и распределительных переменных. Функция \(g(x_t,i_t)\) задаёт эффект процентной ставки. Важная HANK-мотивированная особенность состоит в том, что сила трансмиссии ставки зависит от распределительного блока:
\[
\tau_t
=
\tau_0
\left(
1+\chi_\kappa \overline{\kappa}_t+\chi_\ell \ell_t
\right).
\]
Поэтому одни и те же движения ставки могут иметь разный эффект на инфляционный разрыв и разрыв выпуска в зависимости от распределения домохозяйств.

Центральный банк не наблюдает \(x_t\) напрямую. Он получает два типа шумных сигналов. Агрегатные наблюдения имеют вид
\[
o_t^{agg}
=
\left(
\tilde \pi_t^{gap},
\tilde y_t^{gap}
\right)
=
M_{agg}x_t+\nu_t^{agg}.
\]
В сценариях с распределительной информацией дополнительно доступны шумные распределительные сигналы:
\[
o_t^{dist}
=
\left(
\widetilde{\overline{\kappa}}_t,
\tilde \ell_t
\right)
=
M_{dist}x_t+\nu_t^{dist}.
\]

Информационное множество регулятора:
\[
\mathcal I_t =
\sigma(o_0,\ldots,o_t,i_0,\ldots,i_{t-1}).
\]
Здесь \(o_t\) обозначает тот набор сигналов, который доступен в данном информационном режиме.

Правило политики строится по информационному состоянию
\[
s_t = \phi(\mathcal I_t),
\qquad
i_t = \mu(s_t).
\]

Функция потерь считается по истинным реализованным переменным, а не по их наблюдаемым сигналам:
\[
L_t
=
\lambda_\pi(\pi_t^{gap})^2
+
\lambda_y(y_t^{gap})^2
+
\lambda_i(i_t-i_{t-1})^2.
\]
Накопленные потери правила равны
\[
J(\mu,\phi)
=
\mathbb E
\sum_{t=0}^{T}
\beta^t L_t.
\]
В основной спецификации используется конечный горизонт и локальная задача стабилизации без явной нижней границы процентной ставки. В вычислительной реализации ставка ограничивается симметричным техническим диапазоном, чтобы исключить численно неинтерпретируемые траектории.
